\section{Proof}

\subsection{The source-defined systems and their maps}

For a smooth closed form \(\sigma\) in a pseudoeffective class, let
\(\mathcal C(\sigma)\) denote the prime test-curve space
\(\PSH^{\NA}(X,\sigma)'\) of \cite[Definition~3.2]{XiaOps}.
If \(\sigma'-\sigma\) is K\"ahler, its transition map is
\[
 p_{\sigma,\sigma'}(U)_\tau=P_{\sigma'}[U_\tau]_{\mathcal I},
 \qquad \tau<c_U.
\]
It preserves the endpoint and every multiplier ideal at every fixed
level. The transition identities are part of
\cite[Definition~3.5]{XiaOps}. Thus
\begin{equation}\label{eq:p143n-system}
 \mathcal N(\sigma)=\varprojlim_{\eta\in\mathcal K(X)}
 \mathcal C(\sigma+\eta),
\end{equation}
where \(\mathcal K(X)\) is the set of K\"ahler forms and the order is
reversed: a smaller positive perturbation determines a component at a
larger perturbation. Whenever two perturbations are not comparable, a
sufficiently small positive multiple of a fixed K\"ahler form lies
strictly below both. All inverse limits here are limits of sets with
these specified maps.

For \(\sigma'-\sigma\) K\"ahler, the natural injection
\cite[(3.5)]{XiaOps}
\[
 j_{\sigma,\sigma'}:\mathcal N(\sigma)\longrightarrow
 \mathcal N(\sigma')
\]
is obtained by restriction to the components whose total form is
\(\sigma'+\nu\), with \(\nu\) K\"ahler. In particular,
\begin{equation}\label{eq:p143n-restriction}
 \bigl(j_{\sigma,\sigma'}\Gamma\bigr)^{\sigma'+\nu}
 =\Gamma^{\sigma'+\nu}.
\end{equation}
It preserves the raw function on every semivaluation, not merely on
divisorial ones: each expression in \eqref{eq:p143n-raw} uses exactly
the same multiplier ideals after passage to a common component.

We also use the standard cohomological identification stated in the
introduction of \cite[p.~3, after the definition of the projective
limit]{DXZ23}. If \(\sigma'=\sigma+\ddc h\), with \(h\) smooth,
the normalized level corresponding to \(U_\tau\) in the component
\(\sigma+\eta\) is
\begin{equation}\label{eq:p143n-reference-change}
 P_{\sigma'+\eta}[U_\tau-h]_{\mathcal I}.
\end{equation}
This is the source-defined change of representative of the test curve,
not an assertion that \(U_\tau-h\) itself is normalized
\(\mathcal I\)-model. The multiplier-ideal envelope preserves the
ideals of its argument, and a bounded smooth summand does not change
any multiplier ideal. Indeed, on a relatively compact coordinate
neighbourhood, the factor \(e^{mh}\) is bounded above and below by
positive constants, which does not affect local integrability.
Consequently \eqref{eq:p143n-reference-change} has the same ideals as
\(U_\tau\). The inverse construction adds \(h\) and re-envelopes;
the normalized \(\mathcal I\)-model potential is determined by these
ideals. The same description commutes with transition maps and preserves
the endpoint and raw function. This explains precisely which feature
of the known cohomological identification is needed below.

\subsection{A componentwise proof of inverse-limit flattening}

\begin{lemma}\label{lem:p143n-flatten}
Fix \(\sigma\) in a pseudoeffective class and a K\"ahler form
\(\omega\). The natural map
\[
 \mathcal N(\sigma)\longrightarrow
 \varprojlim_{\varepsilon>0}\mathcal N(\sigma+\varepsilon\omega)
\]
is a bijection. The maps on the right are
\(j_{\sigma+\varepsilon\omega,\sigma+\delta\omega}\) for
\(0<\varepsilon<\delta\). No mass condition on the given components
is required.
\end{lemma}

\begin{proof}
We write out the inverse rather than appealing to ordinary density or
to a formal interchange of two limits. Suppose that
\((G_\varepsilon)_{\varepsilon>0}\) is a compatible family on the
right. Given a K\"ahler form \(\eta\), choose \(\varepsilon>0\)
so small that \(\eta-\varepsilon\omega\) is K\"ahler, and put
\begin{equation}\label{eq:p143n-diagonal}
 H^{\sigma+\eta}:=(G_\varepsilon)^{\sigma+\eta}
 \in\mathcal C(\sigma+\eta).
\end{equation}
Such a choice exists: the least eigenvalue of \(\eta\) relative to
\(\omega\) has a positive minimum on compact \(X\). The expression
on the right is the component of \(G_\varepsilon\) at the positive
additional perturbation \(\eta-\varepsilon\omega\).

To show independence of choice, take two admissible numbers
\(\varepsilon_1,\varepsilon_2\), and choose
\(0<\delta<\min(\varepsilon_1,\varepsilon_2)\). Compatibility and
\eqref{eq:p143n-restriction} give
\[
 (G_{\varepsilon_i})^{\sigma+\eta}
 =\bigl(j_{\sigma+\delta\omega,\sigma+\varepsilon_i\omega}
 G_\delta\bigr)^{\sigma+\eta}
 =(G_\delta)^{\sigma+\eta},\qquad i=1,2.
\]
All displayed components exist, since
\(\eta-\delta\omega\) is K\"ahler. Thus
\eqref{eq:p143n-diagonal} is well-defined.

If \(\eta'-\eta\) is K\"ahler, choose one small
\(\varepsilon\) admissible for \(\eta\), hence also for \(\eta'\).
The internal compatibility of \(G_\varepsilon\) gives
\[
 p_{\sigma+\eta,\sigma+\eta'}(H^{\sigma+\eta})
 =p_{\sigma+\eta,\sigma+\eta'}
       ((G_\varepsilon)^{\sigma+\eta})
 =(G_\varepsilon)^{\sigma+\eta'}=H^{\sigma+\eta'}.
\]
Its endpoints agree as well: the endpoints of two \(G_\varepsilon\)
agree by comparison with a smaller \(G_\delta\), and their internal
components have that same finite endpoint. We have therefore constructed
\(H\in\mathcal N(\sigma)\) by \eqref{eq:p143n-system}.

For each \(\varepsilon>0\) and each additional K\"ahler form
\(\nu\), choose \(0<\delta<\varepsilon\) in the construction of
the total component \(\sigma+\varepsilon\omega+\nu\). Then
\[
 \bigl(j_{\sigma,\sigma+\varepsilon\omega}H\bigr)
            ^{\sigma+\varepsilon\omega+\nu}
 =(G_\delta)^{\sigma+\varepsilon\omega+\nu}
 =(G_\varepsilon)^{\sigma+\varepsilon\omega+\nu}.
\]
Equality of all components proves that re-embedding \(H\) recovers
\(G_\varepsilon\). Conversely, starting with \(H\) and restricting
its components gives back \(H\) in \eqref{eq:p143n-diagonal}.
This proves both inverse identities.
\end{proof}

The proof used the prime spaces, whose definition permits zero mass.
The identification with the positive-mass inverse-system presentation
is precisely \cite[Lemma~3.7]{XiaOps}. Its displayed diagonal is the
component with intermediate form \(\sigma+\eta/2\) and final form
\(\sigma+\eta\). We do not assume that an arbitrary isolated prime
test curve has positive mass, or apply positive-volume recovery at
\(\sigma\). In particular, \cite[Lemma~3.10]{XiaOps} is not an
input to \cref{lem:p143n-flatten}.

\subsection{Passage from the K\"ahler cones to their intersection}

We prove \cref{thm:p143n-nef}. First let
\(\Gamma\in\mathcal N(\theta)\). For each \(\varepsilon>0\), its
image \(j_{\theta,\theta+\varepsilon\omega}\Gamma\) is a full
DXZ object in the K\"ahler class \(\alpha+\varepsilon\beta\).
If its displayed representative is not positive, first use
\eqref{eq:p143n-reference-change} to choose a positive representative.
Neither step changes the raw function. By
\cref{hyp:p143n-kahler},
\[
 F_\Gamma\in\PSH_{\mathrm{MP}}(\alpha+\varepsilon\beta)
 \quad\hbox{for every }\varepsilon>0.
\]
Thus \(F_\Gamma\in\mathcal Q_\omega(\alpha)\).

Conversely, let \(F\in\mathcal Q_\omega(\alpha)\). For every
\(\varepsilon>0\), \cref{hyp:p143n-kahler} supplies a unique
\(G_\varepsilon\in\mathcal N(\theta+\varepsilon\omega)\)
with raw function \(F\), on all of \(X^\beth\). If
\(0<\varepsilon<\delta\), the object
\(j_{\theta+\varepsilon\omega,\theta+\delta\omega}
G_\varepsilon\) still has raw function \(F\). Uniqueness in the
K\"ahler class \(\alpha+\delta\beta\) gives
\begin{equation}\label{eq:p143n-coherence}
 j_{\theta+\varepsilon\omega,\theta+\delta\omega}
 G_\varepsilon=G_\delta.
\end{equation}
All endpoints equal the same finite number \(F(o)\). Apply
\cref{lem:p143n-flatten} to obtain \(\Gamma\in\mathcal N(\theta)\)
whose images are exactly the \(G_\varepsilon\). Using any one of
these images in \eqref{eq:p143n-raw} yields \(F_\Gamma=F\)
at every semivaluation. This is not a passage from equality on a dense
set to equality of arbitrary functions.

Finally the natural injections preserve and reflect the DXZ order
\cite[Definition~4.1 and Lemma~4.2]{XiaOps}. Equivalently, one may
check order componentwise using the same choice of \(\varepsilon\)
in \eqref{eq:p143n-diagonal}. If \(\Gamma\leq\Gamma'\), their
images in every K\"ahler perturbation are ordered, so the conditional
K\"ahler comparison gives \(F_\Gamma\leq F_{\Gamma'}\).
If the latter inequality holds, that comparison gives the order of
the images, and the componentwise construction gives
\(\Gamma\leq\Gamma'\). This proves order reflection, and also
injectivity. Together with the preceding surjectivity it proves the
asserted order isomorphism.

\subsection{Auxiliary-form and cohomological independence}

We first record a direct consequence of the model definition in
\cite[Definitions~4.1.2 and 4.3.1]{MP24}. If \(\gamma,\gamma'\)
are K\"ahler classes and \(\gamma'-\gamma\) is nef, then
\begin{equation}\label{eq:p143n-mp-monotone}
 \PSH_{\mathrm{MP}}(\gamma)\subset
 \PSH_{\mathrm{MP}}(\gamma')
\end{equation}
as sets of actual functions. For a model function represented by a
vertical divisor \(D\) on a dominating test configuration
\(\mathcal X\), the class \(\gamma_{\mathcal X}+D\) is relatively
nef. Adding the pullback of the nef class \(\gamma'-\gamma\)
preserves relative nefness. The function represented by \(D\) has
not changed. The same decreasing net of model functions consequently
proves \eqref{eq:p143n-mp-monotone}. No maximum or minimum of models
is needed.

Let \(\omega_1,\omega_2\) be K\"ahler forms and suppose
\(F\in\mathcal Q_{\omega_1}(\alpha)\). For a given
\(\varepsilon>0\), choose \(\delta>0\) so small that
\(\varepsilon\omega_2-\delta\omega_1\) is K\"ahler. Both
\(\alpha+\delta\{\omega_1\}\) and
\(\alpha+\varepsilon\{\omega_2\}\) are K\"ahler, and their
difference is K\"ahler. Equation \eqref{eq:p143n-mp-monotone} gives
\[
 F\in\PSH_{\mathrm{MP}}(\alpha+\delta\{\omega_1\})
 \subset\PSH_{\mathrm{MP}}(\alpha+\varepsilon\{\omega_2\}).
\]
Since \(\varepsilon\) was arbitrary, this proves one inclusion of
the intersections; reversing the forms proves equality. This also
shows that using all positive real perturbations or the cofinal
sequence \(\varepsilon=1/j\) gives the same intersection.

The MP model criterion depends only on the cohomology classes, so the
intersection depends only on \(\alpha\). If \(\theta'\) is another
smooth representative of \(\alpha\), the \(\ddc\)-lemma on the
compact K\"ahler manifold gives \(\theta'=\theta+\ddc h\) with
\(h\) smooth and real. The DXZ identification
\eqref{eq:p143n-reference-change} preserves the raw function.
Therefore the isomorphisms obtained from \(\theta\) and \(\theta'\)
are the same under that identification. Likewise the inverse
constructed using a different auxiliary form has the same raw
function and hence is the same object by injectivity. This proves
all independence assertions in \cref{thm:p143n-nef}.
